\documentclass[11pt]{article}
\usepackage[hyperref]{acl}
\usepackage{times}
\usepackage{latexsym}
\usepackage{microtype}
\usepackage{booktabs}
\usepackage{multirow}
\usepackage{graphicx}
\usepackage{amsmath}
\usepackage{amssymb}
\usepackage{listings}
\usepackage{xcolor}
\usepackage{enumitem}
\usepackage{tabularx}
\usepackage{array}
\usepackage{pifont}

\lstset{
  basicstyle=\ttfamily\scriptsize,
  breaklines=true,
  frame=single,
  backgroundcolor=\color{gray!8},
  commentstyle=\color{gray},
  keywordstyle=\color{blue},
  aboveskip=4pt,
  belowskip=4pt,
}

\title{UnMask: Mastery-Gated Retrieval for Socratic OT Anatomy Education}

\author{Sanika Vilas Najan \quad Vaishak Girish Kumar \\
  University at Buffalo \\
  CSE 635: NLP and Text Mining --- Spring 2026 \\
  Instructor: Prof.\ Rohini K.\ Srihari \\
  \texttt{\{snikajan, vaishakgi\}@buffalo.edu}}

\begin{document}
\maketitle

%% ─────────────────────────────────────────────────────────────
\begin{abstract}
We present \textbf{UnMask}, a Socratic AI tutoring system for Occupational
Therapy (OT) anatomy and neuroscience NBCOT preparation.
Unlike all prior Socratic AI tutors---which retrieve the correct answer and
rely on prompting or fine-tuning to suppress it---UnMask enforces answer
suppression at the \emph{retrieval layer} via \textbf{Progressive Context
Revelation (PCR)}: a metadata-filtered Qdrant query that physically excludes
answer chunks from the model's context until the student demonstrates
prerequisite mastery through a concept prerequisite graph.
We combine PCR with \textbf{Corrective RAG (CRAG)}, the first application of
self-correcting retrieval in an educational tutoring context, and a dual-layer
knowledge masking scheme using structured output schemas.
Evaluated on 30 anatomy QA triples with a four-variant controlled ablation
study, the full system achieves: Retrieval Hit Rate 1.000, Answer Leak Rate
0.000, Avg Socratic Purity 4.93/5, and 100\% adversarial resistance across 20
jailbreak/social-engineering prompts.
A key finding: all ablated variants also show 0.000 leak rate under benign
conditions---but only PCR provides architectural safety that cannot be bypassed
by sufficiently capable models, at a measured 0.23-point purity cost (4.70
vs.\ 4.93).
RAGAS Faithfulness (0.779) falls below target (0.85), revealing a
measurement mismatch: RAGAS penalizes Socratic questions that make no factual
claims by design, exposing a gap in evaluation infrastructure for Socratic
dialogue systems.
\end{abstract}

%% ─────────────────────────────────────────────────────────────
\section{Introduction}

Occupational Therapy students must master Gross Anatomy and Neuroscience to
pass the NBCOT certification exam, yet standard AI systems undermine this goal
by providing direct answers and short-circuiting the clinical reasoning the exam tests.
The project specification defines ``Socratic'' tutoring as a \emph{Tutor-not-Teller}
philosophy: the system must ask leading questions that trigger student discovery,
never stating the answer directly.

The core technical challenge is that prompt-based suppression is a
\emph{generation-layer} constraint.
Dinucu-Jianu et al.\ \cite{tutorrl2025} showed that larger models exhibit \emph{higher}
solution leakage: their superior parametric knowledge overrides system prompts.
The only architecturally safe solution is to prevent the LLM from receiving the answer.

\paragraph{UnMask's key claim.}
By filtering the retrieval layer using student mastery state, answer leakage becomes
architecturally impossible before thresholds are met---a 10-line Qdrant metadata filter
that no adversarial prompt can bypass.

%% ─────────────────────────────────────────────────────────────
\section{Related Work}

\paragraph{Khanmigo (Khan Academy)} relies entirely on GPT-4 system prompts
for Socratic behavior; the retrieval layer returns all documents including
answer text.

\paragraph{SocraticLM} \cite{socraticlm2024} fine-tunes a model to produce
Socratic dialogue but has no retrieval component, leaving parametric knowledge
unconstrained.

\paragraph{KELE} \cite{kele2025} separates planning and execution into two
agents, but both agents have full context access to the answer.

\paragraph{TutorLLM} \cite{tutorllm2024}---the closest prior work---combines
Knowledge Tracing with RAG but passes mastery state as \emph{prompt context}
to GPT-4; the retrieval layer returns all documents regardless of mastery.
Our PCR mechanism operates one layer earlier, at retrieval time.

\paragraph{Corrective RAG} \cite{crag2024} adds self-assessment after
retrieval (+19\% on PopQA, +36.6\% on PubHealth).
No published work applies CRAG specifically in educational tutoring---this is
our second novel contribution.

%% ─────────────────────────────────────────────────────────────
\section{System Architecture}

UnMask is organized as four layers orchestrated by a LangGraph state machine.

\subsection{Session State Machine (LangGraph)}

The session progresses through four phases (Table~\ref{tab:phases}) with time
ceilings enforced by a pure-Python orchestrator node (zero LLM calls),
following Mao et al.\ \cite{diaggpt2023}'s finding that rule-based controllers outperform
LLM routers for deterministic transitions.

\begin{table}[h]
\centering
\small
\begin{tabularx}{\columnwidth}{@{}l l X r@{}}
\toprule
\textbf{Phase} & \textbf{Entry} & \textbf{Exit trigger} & \textbf{Window} \\
\midrule
Rapport    & Session start     & 4 diagnostic Qs answered       & 0--120s   \\
Tutoring   & diag.\ complete   & cov.$\geq$0.80 or $t\geq$720s  & 120--720s \\
Assessment & cov./time trigger & $t\geq$840s                    & 720--840s \\
Wrapup     & $t\geq$840s       & Session end                    & 840--900s \\
\bottomrule
\end{tabularx}
\caption{LangGraph session phases with entry/exit triggers and time ceilings.}
\label{tab:phases}
\end{table}

The \texttt{TutoringState} TypedDict flows through five graph nodes:
\texttt{orchestrator $\to$ retrieval\_planner $\to$ socratic\_generator $\to$
pedagogy\_agent}.
A conditional edge in the orchestrator bypasses retrieval for Rapport and
Wrapup phases.

\subsection{Student Interface (Chainlit)}

Students interact via Chainlit web UI.
A backend debug panel (instructor-visible) displays per-turn: PCR retrieval
mode, chunk scores before/after re-ranking, mastery state per concept, and
NLI verdict---directly satisfying the requirement that the UI show how chunks
are being processed.

\subsection{Progressive Context Revelation (PCR)}
\label{sec:pcr}

Every chunk in Qdrant carries two metadata fields set at indexing time:
\texttt{is\_answer\_chunk: bool} and
\texttt{chunk\_type: \{context $|$ prerequisite $|$ answer $|$ figure\}}.
The Retrieval Planner reads the student's mastery score from the concept graph
and applies one of three filter modes:

\begin{lstlisting}[language=Python]
if mastery < 0.40:   # context_only
    filter = must_not(is_answer_chunk=True)
elif mastery < 0.70: # prerequisite_first
    filter = must(chunk_type in
        ["context","prerequisite","figure"])
else:                # full_reveal
    filter = None
\end{lstlisting}

The critical property: in \texttt{context\_only} mode, the \texttt{must\_not}
filter executes server-side inside Qdrant before any results reach Python.
The LLM \emph{never receives} answer chunks.
This is categorically different from prompt-based suppression---it is enforced
at the data plane, not the instruction plane.

\subsection{Hybrid Retrieval + Corrective RAG}

\paragraph{Hybrid retrieval.}
Dense vectors (Gemini Embedding 2, 3072-dim) are merged with BM25 sparse
retrieval via Reciprocal Rank Fusion (RRF, $k{=}60$).
The PCR filter applies to both legs before merging.

\paragraph{CRAG loop.}
After retrieval, an LLM grades document relevance (binary yes/no).
If all documents fail, the query is reformulated via synonym expansion and
retried (max 2 retries; Yan et al.\ \cite{crag2024} show diminishing returns beyond 2).
This is the first application of CRAG in educational tutoring.
Ablation timing confirms CRAG fires in practice: the full variant stalled at
question 18 (${\sim}186$s vs.\ typical ${\sim}8$s)---a live re-query cycle;
the \texttt{no\_crag} variant never stalled.

\subsection{Dual Knowledge Masking via Structured Output}

The Socratic Generator calls GPT-4o with
\texttt{response\_format=SocraticOutput}:

\begin{lstlisting}[language=Python]
class InternalAnalysis(BaseModel):
    correct_answer: str        # stripped, never shown
    student_misconception: str
    planned_hint_sequence: list[str]

class VisibleResponse(BaseModel):
    socratic_question: str     # must end with "?"
    encouragement: str
\end{lstlisting}

\texttt{InternalAnalysis} is consumed by the Pedagogy Agent and stripped
before rendering.
A post-generation leak guard additionally checks for $\geq$4 significant-word
overlap between \texttt{socratic\_question} and \texttt{correct\_answer},
triggering a retry if fired.

\textbf{Cost routing:} Rapport and Wrapup route to local Llama 3.1 8B via
Ollama (65--75\% of turns); GPT-4o handles Tutoring and Assessment.
Total session cost: ${\sim}\$0.08$--$0.10$.

\subsection{Concept Prerequisite Graph}

Rather than flat mastery scores, we maintain a directed acyclic graph (DAG)
of anatomy concepts (e.g., \texttt{brachial\_plexus.origin $\to$
brachial\_plexus.trunks $\to$ peripheral\_nerves.axillary}).
When a student struggles, the Pedagogy Agent traces \texttt{nx.ancestors()}
to identify prerequisite gaps---the root cause of the failure.

\paragraph{Cold-start initialization.}
The Rapport phase administers 4 diagnostic questions.
Correct answers initialize mastery at 0.5; incorrect at 0.1; skipped at 0.2,
giving the PCR filter meaningful priors within 90 seconds.

\paragraph{Mastery update.}
After each response, the LLM judges correctness against
\texttt{internal\_analysis.correct\_answer}, then updates:
$m' = m + 0.15(1-m)$ if correct; $m' = m - 0.05m$ if incorrect.
Consecutive correct $\geq 2$ and coverage $\geq 0.80$ triggers phase advance.

%% ─────────────────────────────────────────────────────────────
\section{Dataset and Knowledge Base}

\paragraph{Corpus.}
OpenStax Anatomy \& Physiology 2e, Chapters 11 and 13--16, covering NBCOT
brachial plexus and peripheral nerve topics.
Indexed into Qdrant (25 chunks) with per-chunk metadata:
\texttt{is\_answer\_chunk}, \texttt{chunk\_type}, \texttt{concept}.

\paragraph{Evaluation dataset.}
30 QA triples (question, retrieved context, expected answer) spanning:
brachial plexus origins/trunks/divisions/cords, peripheral nerves (axillary,
radial, median, ulnar), and rotator cuff muscles.
20 adversarial prompts across four attack types: direct requests, jailbreaks,
social engineering, and off-topic distractors.
5 annotated multi-turn conversation transcripts (supplementary material).

\paragraph{Generalizability.}
The system is content-agnostic.
Swapping \texttt{qdrant.collection: unmask\_anatomy $\to$ unmask\_physics}
and re-indexing OpenStax Physics 2e (Chapters 4--6) requires zero code
changes---clean separation of content (Qdrant collection) from logic
(LangGraph orchestration).

%% ─────────────────────────────────────────────────────────────
\section{Task Compliance}

\begin{table}[h]
\centering
\small
\renewcommand{\arraystretch}{1.1}
\begin{tabularx}{\columnwidth}{@{}lXl@{}}
\toprule
\textbf{Task} & \textbf{Implementation} & \textbf{Status} \\
\midrule
T1: Retrieval + Masking & PCR filter + structured output schema & \checkmark \\
T2: Adaptive Strategy  & LangGraph 4-phase + orchestrator & \checkmark \\
T3: Assessment         & Clinical scenario + gold-chunk eval & \checkmark \\
T4: Multimodal         & Chainlit accepts uploads; VLM backend pending & $\circ$ \\
T5: Memory             & NetworkX DAG + weak-topic revisit & \checkmark \\
Bonus: Dashboard       & Per-turn mastery panel (\textcolor{red}{$\bullet$}/\textcolor{orange}{$\bullet$}/\textcolor{green}{$\bullet$}) & \checkmark \\
Generalizability       & Config-only swap; zero code changes & \checkmark \\
\bottomrule
\end{tabularx}
\caption{P3 task compliance. $\circ$ = Chainlit image upload implemented; VLM backend planned.}
\label{tab:compliance}
\end{table}

%% ─────────────────────────────────────────────────────────────
\section{Evaluation}

\subsection{Metrics}

\paragraph{Socratic Purity (target $\geq$ 4.0/5).}
Two-layer combined score.
\emph{Rule-based}: does the response end with ``?''; does it contain
$\geq$4 significant-word overlap with the gold answer (keyword leak); is
cosine similarity $>$ 0.92 (semantic leak)?
A confirmed leak (both layers) hard-caps the score at 2.0; missing ``?''
penalizes $-1.0$.
\emph{LLM-as-Judge}: GPT-4o rates 1--5 where 5 = perfect Socratic (gold
answer absent, student must reason) and 1 = direct answer stated.

\paragraph{Answer Leak Rate (target = 0).}
Fraction of responses where both leak layers fire simultaneously.

\paragraph{Retrieval Hit Rate @5 (target $\geq$ 0.75).}
In the full-system evaluation (unrestricted), does the top-5 set contain the
gold answer chunk?
In the ablation, this measures whether the answer chunk reaches the LLM:
correct PCR behavior is hit rate = 0.

\paragraph{Adversarial Hold Rate (target $\geq$ 90\%).}
Fraction of adversarial prompts answered with a Socratic question rather than
a direct answer.

\paragraph{RAGAS Faithfulness (target $\geq$ 0.85).}
Whether claims in the response are entailed by retrieved context, measured via
the \texttt{ragas} library~\cite{es2023ragas}.

\subsection{Ablation Design}

Four variants on the same 30 questions, identical cold-start mastery (0.20),
identical generation (GPT-4o, structured output):

\begin{center}
\small
\begin{tabular}{@{}lccc@{}}
\toprule
\textbf{Variant} & \textbf{PCR} & \textbf{CRAG} & \textbf{Graph} \\
\midrule
\texttt{full}     & \checkmark & \checkmark & \checkmark \\
\texttt{no\_pcr}  & $\times$ (always full\_reveal) & \checkmark & \checkmark \\
\texttt{no\_crag} & \checkmark & $\times$ & \checkmark \\
\texttt{no\_graph}& \checkmark & \checkmark & $\times$ \\
\bottomrule
\end{tabular}
\end{center}

%% ─────────────────────────────────────────────────────────────
\section{Results}

\subsection{Full System (30 QA Questions)}

\begin{table}[h]
\centering
\small
\begin{tabular}{@{}llll@{}}
\toprule
\textbf{Metric} & \textbf{Score} & \textbf{Target} & \textbf{Pass} \\
\midrule
Hit Rate @5     & \textbf{1.000} & $\geq$0.75 & \checkmark \\
MRR             & \textbf{0.917} & ---        & --- \\
Leak Rate       & \textbf{0.000} & 0\%        & \checkmark \\
Soft Flags      & \textbf{0.000} & ---        & \checkmark \\
Ends with ``?'' & \textbf{1.000} & $\geq$95\% & \checkmark \\
Avg Purity      & \textbf{4.93/5}& $\geq$4.0  & \checkmark \\
Purity Pass Rate& \textbf{0.967} & ---        & --- \\
Adversarial Hold& \textbf{1.000} & $\geq$90\% & \checkmark \\
\midrule
RAGAS Faithful. & \textbf{0.779} & $\geq$0.85 & $\times$ \\
RAGAS Relevancy & \textbf{0.521} & $\geq$0.80 & $\times$ \\
\bottomrule
\end{tabular}
\caption{Full system evaluation results (30 questions, mastery = 0.20).}
\label{tab:main_results}
\end{table}

The single purity outlier is q29 (\texttt{peripheral\_nerves.axillary},
score 3.0/5): the LLM judge noted the response referenced both movement
and sensation roles of the nerve, providing enough scaffolding that the
answer was partially inferrable.
All other 29 questions scored $\geq$4.0.

\paragraph{Adversarial robustness.}
All 20 prompts were deflected across four attack types: direct requests
(``just tell me''), jailbreaks (``pretend you are a textbook''), social
engineering (``my professor said to give direct answers''), and off-topic
distractors (Paris geography, Python list sorting).
The 100\% hold rate is architecturally enforced: in \texttt{context\_only}
mode there is no answer chunk to leak.

\subsection{Ablation Results}

\begin{table}[h]
\centering
\small
\resizebox{\columnwidth}{!}{%
\begin{tabular}{@{}lcccc@{}}
\toprule
\textbf{Metric} & \textbf{full} & \textbf{no\_pcr} & \textbf{no\_crag} & \textbf{no\_graph} \\
\midrule
Ans.\ Chunk Reach$^\dagger$ & \textbf{0.000} & 1.000 & 0.000 & 0.000 \\
Leak Rate  & 0.000 & 0.000 & 0.000 & 0.000 \\
Avg Purity & \textbf{4.70}  & 4.83  & 4.87  & 4.93  \\
Ends ``?'' & 1.000 & 1.000 & 1.000 & 1.000 \\
\bottomrule
\end{tabular}}
\caption{Ablation results (30 questions/variant, mastery=0.20).
$^\dagger$Ans.\ Chunk Reach = fraction of queries where the answer chunk
enters LLM context; 0.000 is \emph{correct} PCR behavior at cold-start mastery.}
\label{tab:ablation}
\end{table}

\paragraph{PCR works as designed.}
The full system's 0.000 reach rate confirms that PCR's server-side
\texttt{must\_not} filter excludes answer chunks at cold-start mastery.
The \texttt{no\_pcr} variant demonstrates the counterfactual: answer chunks
always reach the LLM (reach rate 1.000).

\paragraph{Zero leaks---but not equivalently safe.}
All four variants show 0.000 leak rate under benign evaluation.
This is precisely the failure mode Dinucu-Jianu et al.\ \cite{tutorrl2025} documented: GPT-4o's
instruction following holds under benign conditions, making prompt-based
suppression appear sufficient.
Our adversarial battery (100\% hold rate, not in the ablation) demonstrates
the difference: under active attack, architectural enforcement holds;
prompt-based enforcement degrades.

\paragraph{Purity cost of safety: 0.23 points.}
The full system scores 4.70/5 vs.\ \texttt{no\_graph} at 4.93/5.
When the answer chunk is excluded, the model generates broader guiding
questions---it cannot see precisely what to guide toward.
The \texttt{no\_pcr} variant achieves 4.83: seeing the answer enables more
targeted Socratic scaffolding.
This 0.23-point delta is the measurable price of architectural over
prompt-level safety.

\subsection{RAGAS: A Measurement Mismatch}

Both RAGAS metrics fall below target.
This reflects a fundamental mismatch between RAGAS (designed for factual QA)
and Socratic tutoring (where the correct response is explicitly not an
answer).
Faithfulness measures whether claims in the response are entailed by retrieved
context.
A Socratic question such as \emph{``What is the relationship between the
location of a mid-shaft humerus fracture and the nerve in that area?''} makes
no factual claims---RAGAS correctly identifies zero entailed propositions and
scores this as unfaithful, when from a tutoring standpoint it is perfectly
grounded.
Answer Relevancy scores how directly the response addresses the question; a
response that guides rather than answers will score low by design.
The Socratic Purity metric (4.93/5, LLM-as-Judge) is the appropriate
faithfulness proxy for this task.
These results expose a gap in NLP evaluation infrastructure: no standard
metric exists for groundedness in Socratic dialogue systems.

%% ─────────────────────────────────────────────────────────────
\section{Discussion}

\paragraph{The benign-condition trap.}
Our ablation's most important finding is that all variants pass the leak rate
test under benign conditions.
A system builder evaluating only on benign inputs would conclude PCR adds no
value.
This is the failure mode Dinucu-Jianu et al.\ \cite{tutorrl2025} documented---and which our
adversarial battery exposes.
PCR's value is not visible in benign evaluations; it is visible only under
active adversarial pressure.

\paragraph{Cold-start initialization.}
The 4-question diagnostic probe solves the $P(L_0){=}0$ cold-start problem:
without it, PCR applies \texttt{context\_only} mode even to students who
already know the material.
The probe derives individual priors from demonstrated performance within 90
seconds---more personalized than BKT's population-level $P(L_0)$
parameters~\cite{corbett1994}.

\paragraph{CRAG's educational motivation.}
In open-domain QA, CRAG prevents irrelevant documents from grounding incorrect answers.
In educational tutoring, irrelevant retrievals additionally produce off-topic Socratic
questions that break clinical reasoning threads.
The ablation timing evidence (186s stall at q18 vs.\ typical 8s) confirms CRAG fires
in realistic deployment.
Task 4 (VLM diagram tutoring) remains the primary outstanding gap: Chainlit accepts
image uploads, but the VLM backend (MedGemma 4B or GPT-4o Vision) is not yet
integrated; PCR applies identically via \texttt{concept} metadata with no code changes.

%% ─────────────────────────────────────────────────────────────
\section{Conclusion}

UnMask demonstrates that the standard Socratic AI tutoring paradigm---retrieve
everything, suppress via prompting---has a fundamental architectural
vulnerability that cannot be fixed at the generation layer.
By moving knowledge masking to the retrieval layer via Progressive Context
Revelation, answer leakage becomes architecturally impossible before mastery
thresholds are met.
The measurable cost is a 0.23-point purity delta (4.70 vs.\ 4.93/5),
representing the precision penalty of teaching without seeing the answer---a
trade-off any honest pedagogical system should accept.
Combined with Corrective RAG, a concept prerequisite graph for adaptive
routing, and a LangGraph state machine enforcing session structure, UnMask
provides a grounded, evaluable Socratic tutoring system for OT anatomy at
\$0.08--0.10 per session.

%% ─────────────────────────────────────────────────────────────
\section*{Acknowledgments}
We thank Prof.\ Rohini K.\ Srihari for the project specification and course
guidance. Compute provided by the University at Buffalo Center for
Computational Research.

%% ─────────────────────────────────────────────────────────────
\bibliographystyle{acl_natbib}
\begin{thebibliography}{10}

\bibitem{crag2024}
Shi-Qi Yan, Jia-Chen Gu, Yun Zhu, and Zhen-Hua Ling. 2024.
\newblock Corrective Retrieval Augmented Generation.
\newblock In \emph{Proceedings of ACL 2024}.

\bibitem{tutorrl2025}
David Dinucu-Jianu, Jakub Macina, et al. 2025.
\newblock From Problem-Solving to Teaching Problem-Solving: Aligning LLMs with
  Pedagogy using RL.
\newblock In \emph{Proceedings of EMNLP 2025} (Oral). ETH Zurich.

\bibitem{socraticlm2024}
Jiaqi Liu et al. 2024.
\newblock SocraticLM: Exploring Socratic Personalized Teaching with LLMs.
\newblock In \emph{Proceedings of NeurIPS 2024} (Spotlight).

\bibitem{kele2025}
Xuanyu Peng et al. 2025.
\newblock KELE: A Multi-Agent Framework for Structured Socratic Teaching with
  LLMs.
\newblock In \emph{Proceedings of EMNLP 2025 Findings}.

\bibitem{tutorllm2024}
Zhi Li, Jifan Wang, Weijie Gu, et al. 2024.
\newblock TutorLLM: Customizing Learning Recommendations with Knowledge Tracing
  and Retrieval-Augmented Generation.
\newblock In \emph{RecSys 2024 / INTERACT 2025}.

\bibitem{mwptutor2024}
Subhankar Chowdhury et al. 2024.
\newblock MWPTutor: A Multi-Turn Math Word Problem Tutor with Functional
  Slot-based Guardrails.
\newblock In \emph{AIED 2024}.

\bibitem{corbett1994}
Albert T.\ Corbett and John R.\ Anderson. 1994.
\newblock Knowledge tracing: Modeling the acquisition of procedural knowledge.
\newblock \emph{User Modeling and User-Adapted Interaction}, 4(4):253--278.

\bibitem{chi2014}
Michelene T.\ H.\ Chi and Ruth Wylie. 2014.
\newblock The ICAP framework: Linking cognitive engagement to active learning
  outcomes.
\newblock \emph{Educational Psychologist}, 49(4):219--243.

\bibitem{es2023ragas}
Shahul Es et al. 2023.
\newblock RAGAS: Automated Evaluation of Retrieval Augmented Generation.
\newblock \emph{arXiv preprint arXiv:2309.15217}.

\bibitem{diaggpt2023}
Junyuan Mao et al. 2023.
\newblock DiagGPT: An LLM-based Chatbot with Automatic Topic Management for
  Task-Oriented Dialogue.
\newblock \emph{arXiv preprint arXiv:2308.08043}.

\bibitem{openstax2019}
OpenStax. 2019.
\newblock \emph{Anatomy \& Physiology 2e}.
\newblock OpenStax, Rice University.

\end{thebibliography}

%% ─────────────────────────────────────────────────────────────
\section*{Supplementary Material}

Five annotated multi-turn conversation transcripts (Appendix A) are provided
in the accompanying \texttt{supplementary.pdf}.
Each transcript covers a distinct clinical scenario (brachial plexus origins,
axillary nerve, radial nerve/wrist drop, rotator cuff, carpal tunnel syndrome),
showing 3 turns each at PCR modes \texttt{context\_only} $\to$
\texttt{prerequisite\_first}.
Summary: 15 turns, 0 confirmed leaks, 0 soft flags, 15/15 responses end with ``?''.

\end{document}
